% Additional vector figures for classical models.

\newcommand{\DecisionTreeLifecycleFigure}{%
\begin{figure}[htbp]
\centering
\begin{tikzpicture}[
  split/.style={draw=Blue,rounded corners,fill=PaleBlue,align=center,
    minimum width=16mm,minimum height=7mm,inner xsep=2mm,font=\scriptsize},
  leaf/.style={draw=Teal,rounded corners,fill=PaleTeal,align=center,
    minimum width=15mm,minimum height=7mm,inner xsep=2mm,font=\scriptsize},
  weak/.style={draw=Orange,rounded corners,fill=PaleOrange,align=center,
    minimum width=15mm,minimum height=7mm,inner xsep=2mm,font=\scriptsize},
  edge/.style={-{Latex[length=1.7mm]},semithick,color=Ink!70},
  stage/.style={font=\small\bfseries,color=Navy}
]
% Stage 1: one greedy split.
\node[stage] at (0,2.55) {1. Split};
\node[split] (a0) at (0,1.65) {$x_1\le t_1$};
\node[leaf] (a1) at (-0.95,0.45) {mostly $-$};
\node[weak] (a2) at (0.95,0.45) {mixed};
\draw[edge] (a0)--node[left,font=\tiny]{yes}(a1);
\draw[edge] (a0)--node[right,font=\tiny]{no}(a2);

% Stage 2: recursively split eligible non-pure children.
\node[stage] at (5.25,2.55) {2. Grow recursively};
\node[split] (b0) at (5.25,1.65) {$x_1\le t_1$};
\node[split] (b1) at (4.10,0.45) {$x_2\le t_2$};
\node[weak] (b2) at (6.40,0.45) {$x_3\le t_3$};
\node[leaf] (b11) at (3.45,-0.75) {$-$};
\node[leaf] (b12) at (4.75,-0.75) {$+$};
\node[weak] (b21) at (5.75,-0.75) {$-$};
\node[weak] (b22) at (7.05,-0.75) {$+$};
\draw[edge] (b0)--(b1); \draw[edge] (b0)--(b2);
\draw[edge] (b1)--(b11); \draw[edge] (b1)--(b12);
\draw[edge] (b2)--(b21); \draw[edge] (b2)--(b22);
\node[draw=Orange,dashed,rounded corners,fit=(b2)(b21)(b22),
  inner sep=2mm] (weaksubtree) {};

% Stage 3: collapse a weak subtree after validation.
\node[stage] at (11.25,2.55) {3. Prune};
\node[split] (c0) at (11.25,1.65) {$x_1\le t_1$};
\node[split] (c1) at (10.20,0.45) {$x_2\le t_2$};
\node[weak] (c2) at (12.30,0.45) {leaf $+$};
\node[leaf] (c11) at (9.55,-0.75) {$-$};
\node[leaf] (c12) at (10.85,-0.75) {$+$};
\draw[edge] (c0)--(c1); \draw[edge] (c0)--(c2);
\draw[edge] (c1)--(c11); \draw[edge] (c1)--(c12);

\draw[-{Latex[length=2mm]},thick,Blue]
  (1.55,1.25)--node[above,font=\scriptsize,color=Ink]{reduce impurity}(3.10,1.25);
\draw[-{Latex[length=2mm]},thick,Orange]
  (7.75,1.25)--node[above,font=\scriptsize,color=Ink]{validate complexity}(9.35,1.25);
\end{tikzpicture}
\caption{A tree chooses a locally useful split, recursively splits eligible non-pure children while the pre-pruning conditions permit further growth, and may then prune a branch whose added complexity is not supported by validation data. Pruning replaces the highlighted subtree by one leaf.}
\label{fig:tree-growth-pruning}
\end{figure}}

\newcommand{\EnsembleWorkflowFigure}{%
\begin{figure}[htbp]
\centering
\begin{tikzpicture}[
  source/.style={draw=Navy,rounded corners,fill=PaleBlue,align=center,
    minimum width=16mm,minimum height=8mm,font=\scriptsize},
  bag/.style={draw=Blue,rounded corners,fill=PaleBlue,align=center,
    minimum width=18mm,minimum height=7mm,font=\scriptsize},
  boost/.style={draw=Orange,rounded corners,fill=PaleOrange,align=center,
    minimum width=18mm,minimum height=7mm,font=\scriptsize},
  output/.style={draw=Teal,rounded corners,fill=PaleTeal,align=center,
    minimum width=20mm,minimum height=8mm,font=\scriptsize},
  arr/.style={-{Latex[length=1.8mm]},semithick,color=Ink!70},
  lane/.style={font=\small\bfseries,color=Navy,anchor=west}
]
% Bagging and random-forest lane.
\node[lane] at (-0.15,3.00) {Bagging / random forest};
\node[source] (d1) at (0.75,1.55) {training\\data};
\node[bag] (s1) at (3.20,2.25) {bootstrap $B_1$};
\node[bag] (s2) at (3.20,1.55) {bootstrap $B_2$};
\node[bag] (sm) at (3.20,0.85) {bootstrap $B_M$};
\node[bag] (t1) at (6.10,2.25) {tree $T_1$};
\node[bag] (t2) at (6.10,1.55) {tree $T_2$};
\node[bag] (tm) at (6.10,0.85) {tree $T_M$};
\node[output] (agg) at (9.10,1.55) {average\\or vote};
\draw[arr] (d1)--(s1); \draw[arr] (d1)--(s2); \draw[arr] (d1)--(sm);
\draw[arr] (s1)--(t1); \draw[arr] (s2)--(t2); \draw[arr] (sm)--(tm);
\draw[arr] (t1)--(agg); \draw[arr] (t2)--(agg); \draw[arr] (tm)--(agg);
\node[font=\tiny,color=Muted] at (6.10,0.20)
  {RF: sample candidate features at each split};
\node[font=\scriptsize,color=Teal,anchor=west,align=left] at (10.35,1.55)
  {parallel\\variance reduction};

% Boosting lane.
\node[lane] at (-0.15,-0.35) {Gradient boosting};
\node[source] (d2) at (0.75,-1.55) {training\\data};
\node[boost] (h1) at (3.00,-1.55) {fit $h_1$};
\node[boost] (r2) at (5.15,-1.55) {gradients\\from $F_1$};
\node[boost] (h2) at (7.30,-1.55) {fit $h_2$};
\node[boost] (more) at (9.35,-1.55) {$\cdots$};
\node[output] (sum) at (11.45,-1.55) {weighted\\sum};
\draw[arr] (d2)--(h1); \draw[arr] (h1)--(r2); \draw[arr] (r2)--(h2);
\draw[arr] (h2)--(more); \draw[arr] (more)--(sum);
\node[font=\scriptsize,color=Orange,anchor=west,align=left] at (12.65,-1.55)
  {sequential\\error correction};
\end{tikzpicture}
\caption{Bagging fits learners independently on bootstrap samples and aggregates their predictions; a random forest additionally randomises candidate features at each split. Boosting is sequential: each learner follows the loss gradients of the current additive model.}
\label{fig:ensemble-workflows}
\end{figure}}


\section{Learning Problems and a Leakage-Free Workflow}
\label{sec:risk-workflow}

Supervised learning begins with labelled observations
\[
\mathcal D=\{(x_i,y_i)\}_{i=1}^{n},
\qquad x_i\in\mathbb R^d.
\]
The vector $x_i$ records the $d$ features of example $i$. In regression,
$y_i$ is usually real-valued. In classification, $y_i$ belongs to a finite set;
binary labels are commonly encoded as either $\{0,1\}$ or $\{-1,+1\}$.
Each encoding will be stated when it is used.

A model $f_\theta$ is determined by parameters $\theta$ and produces the
prediction $\hat y=f_\theta(x)$. A loss $\ell(y,\hat y)$ measures the error on
one example. Training usually minimises the empirical risk
\[
\widehat R(\theta)
=\frac{1}{n}\sum_{i=1}^{n}\ell\bigl(y_i,f_\theta(x_i)\bigr).
\]
The actual objective is low error on new examples from the same data-generating
distribution. If $(X,Y)$ denotes such an example, the relevant quantity is the
population risk
\[
R(\theta)=\mathbb E_{(X,Y)}
\left[\ell\bigl(Y,f_\theta(X)\bigr)\right].
\]
The population distribution is unknown, so $R(\theta)$ must be estimated using
data that did not determine the fitted model.

\subsection{Training, validation, and test sets}

A sound experiment splits the data before performing any operation that learns
from the data. The three subsets serve different purposes.
\begin{enumerate}
    \item The training set estimates model parameters, such as linear weights,
    tree split points, or neural-network weights.
    \item The validation set selects model classes and hyperparameters, such as
    the number of neighbours, tree depth, or the SVM penalty $C$.
    \item The test set is used once, after all choices have been fixed, to
    estimate the generalisation error of the final procedure.
\end{enumerate}

The word \emph{model} includes the complete preprocessing pipeline. Imputation,
standardisation, feature selection, dimensionality reduction, and resampling
may all learn properties of a dataset. They must therefore be fitted using only
the training data. For example, standardisation uses
\[
z_{ij}=\frac{x_{ij}-\mu_j}{s_j},
\]
where $\mu_j$ and $s_j$ are estimated from the training set. The same values
are then applied to validation and test examples. Estimating them from the full
dataset allows information from the test set to enter training and produces an
optimistic test error.

After hyperparameters have been selected, the final model may be refitted on
the union of the training and validation sets. Every preprocessing parameter
must then be re-estimated on that union, while the test set remains untouched.
Using test results to revise the model turns the test set into another
validation set and invalidates its role as an independent evaluation.

\subsection{The experimental data flow}

For a fixed family of candidate procedures, the data flow is
\[
\mathcal D_{\mathrm{train}}
\xrightarrow{\text{fit preprocessing and model}}
f_{\theta}
\xrightarrow{\mathcal D_{\mathrm{val}}}
\text{select hyperparameters}
\xrightarrow{\text{refit}}
f_{\mathrm{final}}
\xrightarrow{\mathcal D_{\mathrm{test}}}
\text{report error}.
\]
When the dataset is small, cross-validation within the training data can
replace a single validation split. Preprocessing must be refitted inside every
fold using only that fold's training portion.

\TrainingPipelineFigure


\section{Distance, Scaling, and Nearest Neighbours}

K-nearest neighbours (KNN) assumes that nearby observations tend to have
similar labels. Whether this assumption is useful depends first on the feature
representation and distance function. Different distances encode different
notions of similarity and are not interchangeable.

\subsection{Distances and similarities}

For $x,z\in\mathbb R^d$, the Euclidean and Manhattan distances are
\[
d_2(x,z)=\left(\sum_{j=1}^{d}(x_j-z_j)^2\right)^{1/2},
\qquad
d_1(x,z)=\sum_{j=1}^{d}|x_j-z_j|.
\]
Euclidean distance is the straight-line distance and gives relatively large
weight to a large discrepancy in one coordinate. Manhattan distance adds
absolute coordinate-wise discrepancies and is less dominated by such a
discrepancy.

When direction matters more than magnitude, cosine similarity is often more
appropriate for non-zero vectors $x$ and $z$:
\[
\operatorname{cos}(x,z)=\frac{x^\top z}
{\lVert x\rVert_2\lVert z\rVert_2}.
\]
The associated cosine distance is
$d_{\cos}(x,z)=1-\operatorname{cos}(x,z)$. Cosine similarity is not Euclidean
distance. After both vectors have been normalised to unit length, however,
\[
\left\lVert\frac{x}{\lVert x\rVert_2}-
\frac{z}{\lVert z\rVert_2}\right\rVert_2^2
=2\bigl(1-\operatorname{cos}(x,z)\bigr).
\]

For sets or binary features, Jaccard similarity ignores features that are absent
from both observations. If $A$ and $B$ are the sets of non-zero features, then
\[
J(A,B)=\frac{|A\cap B|}{|A\cup B|},
\qquad
d_J(A,B)=1-J(A,B).
\]
This definition is well suited to sparse presence--absence data because a
shared zero does not increase similarity. If $A=B=\varnothing$, the displayed
ratio is undefined; applications that permit empty sets must state a convention,
commonly $J(\varnothing,\varnothing)=1$.

\subsection{Feature scaling}

Distances are sensitive to units. If one feature is measured in thousands and
another lies between $0$ and $1$, the first feature can dominate every
Euclidean distance. Two common transformations are standardisation
\[
z_{ij}=\frac{x_{ij}-\mu_j}{s_j}
\]
and min--max scaling
\[
z_{ij}=\frac{x_{ij}-a_j}{b_j-a_j},
\]
where $a_j$ and $b_j$ are the minimum and maximum of feature $j$ in the
training set. Standardisation does not confine values to a fixed interval, but
it is usually less dependent on two observed endpoints. In either case, all
statistics are estimated from the training data alone. A constant feature has
$s_j=0$ and $b_j-a_j=0$; remove it or map it to a fixed value rather than divide
by zero.

Scaling is not an automatic requirement. Binary indicators, variables with a
meaningful physical scale, and inputs to decision trees need not receive the
same treatment. The choice follows from how the model uses each feature.

\subsection{The curse of dimensionality}
\label{sec:knn-dimensionality}

With a fixed number of observations, the sample becomes sparse as the dimension
increases. In a unit hypercube, a smaller hypercube of side length $r$ occupies
the fraction $r^d$ of the total volume. To enclose a fixed fraction $q$ of the
volume, its side length must be
\[
r=q^{1/d}.
\]
As $d$ increases, $r$ approaches $1$. A neighbourhood must therefore extend
through much of every coordinate range before it contains a fixed fraction of
the observations; it is no longer genuinely local.

Many high-dimensional distributions also exhibit distance concentration: the
relative difference between the nearest and farthest distances shrinks. The
distances remain computable, but their ranking carries less information. For
KNN, high dimensionality primarily causes sparse neighbourhoods and distance
concentration, which can in turn increase variance. More observations, feature
selection, dimensionality reduction, or a structured similarity measure may be
needed.

\subsection{K-nearest neighbours}

For a query $x$, KNN finds the $k$ training observations with the smallest
distance and denotes their index set by $N_k(x)$. A classifier uses a majority
vote,
\[
\hat y
=\operatorname*{arg\,max}_{c}
\sum_{i\in N_k(x)}\mathbf 1(y_i=c),
\]
whereas a regressor commonly uses the mean
\[
\hat y=\frac{1}{k}\sum_{i\in N_k(x)}y_i.
\]
Closer observations can be given more weight. One possible choice is
\[
w_i(x)=\frac{1}{d(x,x_i)+\varepsilon},
\qquad
\hat y
=\frac{\sum_{i\in N_k(x)}w_i(x)y_i}
{\sum_{i\in N_k(x)}w_i(x)}.
\]
where $\varepsilon>0$. The same weights can be used for a weighted class vote.

\KnnGeometryFigure

The value of $k$ controls smoothness. A small $k$ preserves local structure but
is sensitive to noise and individual outliers. A large $k$ is more stable but
may erase a genuine boundary. The rule $k=\sqrt n$ is at most an initial
heuristic; validation or cross-validation must determine the final value. An
odd $k$ reduces ties in binary classification, although ties can still occur
with several classes or repeated distances and therefore require an explicit
tie-breaking rule.

KNN has little explicit training cost but stores the training set. A brute-force
implementation computes $n$ distances in $d$ dimensions for each query, giving
time $O(nd)$ and storage $O(nd)$. Tree-based indices can accelerate low-dimensional
queries, but their benefit often declines in high dimensions.


\section{Decision Trees and Ensembles}

A decision tree partitions the input space recursively. Each internal node asks
a question about a feature, and each leaf stores a prediction. For a continuous
feature, a common split is
\[
x_j\le t,
\]
where $j$ selects a feature and $t$ is a threshold. For the observations $S$ at
the current node, the split produces
\[
S_L=\{(x_i,y_i)\in S:x_{ij}\le t\},
\qquad
S_R=S\setminus S_L.
\]
The two subsets need not have equal size. Training selects the split that most
reduces label impurity or numerical dispersion and repeats the process at the
child nodes.

\subsection{Impurity in classification trees}

Let $p_c$ be the proportion of class $c$ at node $S$. Entropy and Gini
impurity are
\[
H(S)=-\sum_{c=1}^{C}p_c\log_2 p_c,
\qquad
G(S)=1-\sum_{c=1}^{C}p_c^2,
\]
with the convention $0\log 0=0$. Both are zero at a pure node and increase as
the class distribution becomes more mixed.

Suppose a split produces child nodes of sizes $n_L$ and $n_R$, where
$n_S=n_L+n_R$. For an impurity measure $I$, the reduction is
\[
\Delta I
=I(S)-\frac{n_L}{n_S}I(S_L)
-\frac{n_R}{n_S}I(S_R).
\]
When $I=H$, this reduction is the information gain. A greedy tree normally
chooses the feature and threshold with the largest $\Delta I$.

For example, a node containing six positive and four negative examples has
entropy
\[
H(S)=-0.6\log_2 0.6-0.4\log_2 0.4\approx0.971.
\]
Suppose a split sends four positive examples to the left and the remaining two
positive and four negative examples to the right. The weighted child entropy is
\[
\frac{4}{10}\cdot0+
\frac{6}{10}\left(-\frac13\log_2\frac13-
\frac23\log_2\frac23\right)
\approx0.551,
\]
so the information gain is approximately $0.420$. The weighting accounts for
both child purity and child size; inspecting only one child is insufficient.

A classification leaf can output its majority class or estimated class
proportions,
\[
\widehat P(Y=c\mid x)=\frac{n_c}{n_{\mathrm{leaf}}}.
\]
The latter permits threshold adjustment, although probabilities estimated from
small leaves are unstable.

\subsection{Regression trees}

A regression tree uses the same partitioning mechanism. A leaf predicts its
mean response
\[
\bar y_S=\frac{1}{|S|}\sum_{i\in S}y_i.
\]
A common node impurity is the within-node sum of squared errors
\[
I_{\mathrm{reg}}(S)=
\sum_{i\in S}(y_i-\bar y_S)^2.
\]
Training selects the split that maximises
\[
I_{\mathrm{reg}}(S)
-I_{\mathrm{reg}}(S_L)
-I_{\mathrm{reg}}(S_R).
\]
The fitted function is constant within each leaf region, so a single regression
tree produces a piecewise-constant predictor.

\subsection{Stopping and pruning}

If a tree is allowed to keep splitting, a leaf may contain only one observation.
The resulting tree has low training error and high variance. Pre-pruning limits
growth directly through a maximum depth, a minimum leaf size, or a minimum
impurity reduction. Post-pruning first grows a larger tree and then removes
branches whose improvement is not supported by validation data.

\DecisionTreeLifecycleFigure

Cost-complexity pruning penalises the number of leaves:
\[
R_\alpha(T)=R(T)+\alpha |\mathcal L(T)|,
\]
where $R(T)$ is the training loss of tree $T$, $\mathcal L(T)$ is its set of
leaves, and $\alpha\ge0$. Larger values of $\alpha$ favour smaller trees. The
penalty is selected using validation data.

A prediction visits only the path from the root to a leaf, so its cost is
$O(h)$ for tree depth $h$. A balanced tree can have logarithmic depth, whereas
a degenerate tree can have depth linear in its number of leaves. Training cost
depends on how candidate thresholds are generated and sorted; it cannot be
described by a universal $O(dn^2\log n)$ expression.

\subsection{Bagging and random forests}

A deep tree can fit a complicated boundary, but small changes in the training
sample may alter its upper splits. Bagging constructs $M$ bootstrap samples.
Each contains $n$ observations drawn with replacement from the original sample.
The base learners are fitted separately, then averaged in regression or combined
by majority vote or mean class probability in classification:
\[
\bar f(x)=\frac1M\sum_{m=1}^{M}f_m(x).
\]

Suppose the prediction of each base learner at a fixed $x$ has variance
$\sigma^2$, and any pair has approximately the same correlation $\rho$. Then
\begin{equation}
\operatorname{Var}[\bar f(x)]
=\sigma^2\left(\rho+\frac{1-\rho}{M}\right).
\label{eq:correlated-ensemble-variance}
\end{equation}
The variance becomes $\sigma^2/M$ only when the predictors are independent
($\rho=0$). Trees trained on related bootstrap samples remain correlated, so
the reduction is smaller and depends on how effectively the procedure
decorrelates them.

A random forest adds feature randomisation. At each node, the tree searches for
a split among only $m_{\mathrm{try}}$ randomly selected features. This prevents
every tree from repeatedly relying on the same strong feature. The setting
$m_{\mathrm{try}}\approx\sqrt d$ is a common classification default, not a
theoretical optimum; regression implementations often use a different default.
The feature count, tree depth, and leaf size remain hyperparameters.

\EnsembleWorkflowFigure

\subsection{Boosting and gradient boosting}

Bagged learners can be trained in parallel. Boosting adds learners sequentially,
so that each new learner addresses errors of the current ensemble. An additive
model has the form
\[
F_M(x)=F_0(x)+\sum_{m=1}^{M}\eta\,\gamma_m h_m(x),
\]
where $h_m$ is often a shallow tree and $\eta\in(0,1]$ is the learning rate.

Gradient boosting applies this construction to a differentiable loss
$\ell(y,F(x))$. At iteration $m$, it computes the negative gradient at each
training observation,
\[
r_i^{(m)}=-\left.
\frac{\partial\ell(y_i,F(x_i))}{\partial F(x_i)}
\right|_{F=F_{m-1}},
\]
and fits $h_m$ to $\{(x_i,r_i^{(m)})\}$. A line search may choose
\[
\gamma_m=
\operatorname*{arg\,min}_{\gamma}
\sum_{i=1}^{n}
\ell\bigl(y_i,F_{m-1}(x_i)+\gamma h_m(x_i)\bigr),
\]
followed by
\[
F_m(x)=F_{m-1}(x)+\eta\gamma_mh_m(x).
\]
For squared loss $\ell(y,F)=\tfrac12(y-F)^2$, the negative gradient is the
residual $y_i-F_{m-1}(x_i)$. Fitting residuals is therefore the squared-loss
instance of gradient boosting, not its definition for every loss.

Boosting typically reduces bias, while tree depth, learning rate, and the number
of rounds determine capacity and variance. Select the number of rounds by
validation or early stopping.


\section{Regression and Classification from a Linear Score}

Several linear models share the score
\[
s(x)=w^\top x+b.
\]
They differ in how they interpret the score, which loss they minimise, and
whether they impose a margin. Linear regression uses $s(x)$ as a numerical
prediction. Logistic regression maps it to a probability. The Perceptron and
SVM classify according to its sign.

\subsection{Linear regression and the normal equations}
\label{sec:linear-score-regression}

Linear regression takes the score itself as the prediction,
$\hat y_i=s(x_i)$. With the intercept absorbed into an augmented design matrix
$X$ and parameter vector $\theta$, its squared-error objective is
\[
L(\theta)=\frac{1}{2n}\lVert X\theta-y\rVert_2^2.
\]
The derivation of the normal equations, the rank condition for the inverse
formula, and the regularised solution are given in
Section~\ref{sec:least-squares}. Linear regression uses the score directly;
unlike the classifiers below, it neither thresholds the score nor maps it to a
probability.

\subsection{Logistic regression}

Let $y_i\in\{0,1\}$. Logistic regression maps the linear score to a probability:
\[
p_i=P(Y=1\mid x_i)=\sigma(s_i),
\qquad
s_i=w^\top x_i+b,
\qquad
\sigma(z)=\frac{1}{1+e^{-z}}.
\]
Consequently,
\[
P(Y=y_i\mid x_i)=p_i^{y_i}(1-p_i)^{1-y_i}.
\]
Under conditional independence of the observations, maximum conditional
likelihood is equivalent to minimising the mean negative log-likelihood
\begin{align*}
L(w,b)
&=-\frac1n\sum_{i=1}^{n}
\left[y_i\log p_i+(1-y_i)\log(1-p_i)\right]\\
&=\frac1n\sum_{i=1}^{n}
\left[\log(1+e^{s_i})-y_is_i\right].
\end{align*}
This is binary cross-entropy. With the augmented design matrix $X$ and
$p=(p_1,\ldots,p_n)^\top$, its gradient is
\[
\nabla_\theta L=\frac1nX^\top(p-y).
\]
The optimisation procedure is developed in Section~\ref{sec:least-squares} and
the gradient-descent section that follows it.

The threshold $p_i\ge0.5$ is equivalent to $s_i\ge0$, so the decision boundary
is $w^\top x+b=0$. Changing the probability threshold changes the classifier
without refitting the probability model. With labels $y_i\in\{-1,+1\}$, the
same loss is
\[
\ell_{\mathrm{log}}(y_i,s_i)
=\log\bigl(1+e^{-y_is_i}\bigr).
\]

\subsection{The Perceptron}
\label{sec:perceptron}

The Perceptron uses labels $y_i\in\{-1,+1\}$ and predicts
\[
\widehat y_i=\operatorname{sign}(w^\top x_i+b).
\]
An observation is correctly classified when $y_is_i>0$. On a misclassified
observation, the update is
\[
w\leftarrow w+\eta y_ix_i,
\qquad
b\leftarrow b+\eta y_i.
\]
This update is a subgradient step on the Perceptron loss
\[
\ell_{\mathrm{perc}}(y_i,s_i)=\max(0,-y_is_i).
\]
The loss becomes zero as soon as the sign is correct and does not reward a wider
margin. For linearly separable training data, the Perceptron reaches a separating
hyperplane after finitely many updates. It need not find the maximum-margin
hyperplane and does not directly produce class probabilities. On non-separable
data, the original algorithm may continue updating without convergence.

\subsection{Maximum margin and hard-margin SVM}

Many separating hyperplanes may exist for linearly separable data. SVM fixes a
functional-margin scale and selects the hyperplane with the largest geometric
margin. The hard-margin primal problem is
\begin{align*}
\min_{w,b}\quad &\frac12\lVert w\rVert_2^2\\
\text{s.t.}\quad &y_i(w^\top x_i+b)\ge1,
\qquad i=1,\ldots,n.
\end{align*}
The signed distance from $x_i$ to $w^\top x+b=0$ is
$(w^\top x_i+b)/\lVert w\rVert_2$. Under the constraints, the closest training
points lie on $w^\top x+b=\pm1$, and the distance between these support
hyperplanes is
\[
\frac{2}{\lVert w\rVert_2}.
\]
Minimising $\tfrac12\lVert w\rVert_2^2$ therefore maximises the margin. The
hard-margin problem is feasible only when the training data are linearly
separable.

\LinearMarginFigure

\subsection{Soft-margin SVM and hinge loss}

Slack variables $\xi_i\ge0$ allow margin violations, with $C>0$ controlling
their penalty:
\begin{align*}
\min_{w,b,\xi}\quad
&\frac12\lVert w\rVert_2^2+C\sum_{i=1}^{n}\xi_i\\
\text{s.t.}\quad
&y_i(w^\top x_i+b)\ge1-\xi_i,
\qquad \xi_i\ge0.
\end{align*}
For fixed $w,b$, the optimal slack is
\[
\xi_i=\max\bigl(0,1-y_i(w^\top x_i+b)\bigr).
\]
The primal problem is therefore equivalent to
\[
\min_{w,b}\quad
\frac12\lVert w\rVert_2^2
+C\sum_{i=1}^{n}
\underbrace{\max(0,1-y_is_i)}_{\text{hinge loss}}.
\]
A large $C$ penalises violations heavily and favours fitting more training
observations. A small $C$ permits more violations in exchange for a wider
margin. Another common convention writes the objective as
$\lambda\lVert w\rVert_2^2/2+(1/n)\sum_i\ell_i$. The two conventions correspond
under a change of scale; values of $C$ and $\lambda$ cannot be compared without
specifying the objective.

Logistic regression, the Perceptron, and SVM all depend on the signed margin
$ys$, but use different losses:
\[
\begin{aligned}
\ell_{\mathrm{log}}(y,s)&=\log(1+e^{-ys}),\\
\ell_{\mathrm{perc}}(y,s)&=\max(0,-ys),\\
\ell_{\mathrm{hinge}}(y,s)&=\max(0,1-ys).
\end{aligned}
\]
Logistic loss is smooth and admits a probability interpretation. Perceptron
loss vanishes once the sign is correct. Hinge loss requires a functional margin
of at least $1$.

\subsection{The SVM dual}

The dual form of soft-margin SVM depends on the inputs only through inner
products. Introduce multipliers $\alpha_i\ge0$ for the margin constraints and
$\mu_i\ge0$ for $\xi_i\ge0$. The Lagrangian is
\begin{align*}
\mathcal L(w,b,\xi,\alpha,\mu)
=&\ \frac12\lVert w\rVert_2^2+C\sum_i\xi_i\\
&-\sum_i\alpha_i
\left[y_i(w^\top x_i+b)-1+\xi_i\right]
-\sum_i\mu_i\xi_i.
\end{align*}
Stationarity with respect to the primal variables gives
\[
\frac{\partial\mathcal L}{\partial w}=0
\ \Longrightarrow
w=\sum_i\alpha_i y_ix_i,
\]
\[
\frac{\partial\mathcal L}{\partial b}=0
\ \Longrightarrow
\sum_i\alpha_i y_i=0,
\qquad
\frac{\partial\mathcal L}{\partial\xi_i}=0
\ \Longrightarrow
C-\alpha_i-\mu_i=0.
\]
Since $\mu_i\ge0$, the last relation yields $0\le\alpha_i\le C$. Substitution
into the Lagrangian gives the dual problem
\begin{align*}
\max_{\alpha}\quad
&\sum_{i=1}^{n}\alpha_i
-\frac12\sum_{i=1}^{n}\sum_{j=1}^{n}
\alpha_i\alpha_jy_iy_jx_i^\top x_j\\
\text{s.t.}\quad
&0\le\alpha_i\le C,
\qquad
\sum_{i=1}^{n}\alpha_i y_i=0.
\end{align*}
The hard-margin dual has the same objective and equality constraint, but only
$\alpha_i\ge0$ and no upper bound $C$.

Observations with $\alpha_i>0$ appear in
$w=\sum_i\alpha_i y_ix_i$ and are the support vectors. If
$0<\alpha_i<C$, complementary slackness gives
\[
y_i(w^\top x_i+b)=1.
\]
Hence any such support vector gives
\[
b=y_i-\sum_j\alpha_jy_jx_j^\top x_i.
\]
Numerical implementations commonly average this value over several support
vectors satisfying $0<\alpha_i<C$.

\subsection{Kernel SVM}

Both the dual objective and the prediction function require only inner
products. If a feature map $\phi$ exists, a kernel
\[
K(x,z)=\phi(x)^\top\phi(z)
\]
can replace the original inner product without explicitly constructing the
possibly high-dimensional vector $\phi(x)$. The dual becomes
\begin{align*}
\max_{\alpha}\quad
&\sum_i\alpha_i
-\frac12\sum_i\sum_j
\alpha_i\alpha_jy_iy_jK(x_i,x_j)\\
\text{s.t.}\quad
&0\le\alpha_i\le C,
\qquad \sum_i\alpha_i y_i=0,
\end{align*}
and the prediction is
\[
\widehat y(x)=\operatorname{sign}
\left(\sum_{i=1}^{n}\alpha_i y_iK(x_i,x)+b\right).
\]
Only support vectors with $\alpha_i>0$ contribute to this sum.

Common kernels include the linear kernel
\[
K_{\mathrm{linear}}(x,z)=x^\top z,
\]
the polynomial kernel
\[
K_{\mathrm{poly}}(x,z)=(\gamma x^\top z+r)^q,
\qquad \gamma>0,\quad r\ge0,\quad q\in\{1,2,\ldots\},
\]
and the radial basis function kernel
\[
K_{\mathrm{RBF}}(x,z)=
\exp\bigl(-\gamma\lVert x-z\rVert_2^2\bigr),
\qquad \gamma>0.
\]
For a valid kernel, the Gram matrix $K_{ij}=K(x_i,x_j)$ must be positive
semidefinite for every finite set of inputs. In the RBF kernel, a larger $\gamma$ gives each observation a
narrower region of influence and usually permits a more intricate boundary; a
smaller $\gamma$ gives a smoother boundary. The values of $C$ and $\gamma$
jointly control model capacity and should be selected by the leakage-free
validation procedure described earlier.
